
We opened with Thor's observation that hybrid LLM+hammer systems produce 8.2\% unique solutions where neither component succeeds alone, and asked why LLMs and automated provers succeed in different cases. Our mechanistic attribution provides the answer: LLMs exploit natural language hints (29.5 percentage point contribution, ~60\% of the gap) that automated provers ignore, while automated provers provide deterministic formal reasoning without linguistic dependency. This complementarity explains hybrid synergy---LLMs parse NL-rich problems, automated provers solve formal-only goals.

Our controlled ablation study establishes three mechanistic findings: (1) natural language understanding is the dominant mechanism ($\Delta =29.51$\% [20.90\%, 38.11\%], p<$10^{-9}$, large effect Cohen's $h\approx 0.62)$, validating ~60\% contribution, (2) proof depth mechanism provisionally rejected ($\Delta =3.7$\% < 5\% threshold) pending real miniF2F revalidation, (3) tactic budget control validated (CV=0.36) as stable fairness metric.

This mechanistic understanding shifts design priorities for both LLM-guided and hybrid systems. Rather than scaling model architecture or generic pretraining, future work should:

\begin{enumerate}
\item Invest in NL-formal integration (60\% ROI): systematic NL annotation in proof libraries, lightweight $NL\rightarrow tactic$ models for automated provers, semantic tactic suggestion guided by linguistic patterns.
\end{enumerate}

\begin{enumerate}
\item Route by NL content in hybrids: assign NL-rich problems to LLMs (exploit validated 60\% advantage), assign formal-only goals to automated provers.
\end{enumerate}

\begin{enumerate}
\item Revalidate depth hypothesis on real miniF2F: provisional rejection (3.7\% < 5\%) based on mock data (96.3\% shallow-solvable).
\end{enumerate}

\begin{enumerate}
\item Investigate residual 40\% gap: test syntax pattern matching, semantic search, learned heuristics via additional controlled ablations.
\end{enumerate}

\textbf{Broader vision.} Understanding these mechanisms opens the path to NL-aware automated provers that parse hints without LLM costs, hybrid systems with mechanistically-grounded routing, corpus engineering targeting validated mechanisms, and evaluation methodology with gate thresholds enforcing rigor.

Future work should complete real miniF2F validation (H-M1/M2/M3 rerun), test semantic vs syntactic NL understanding via paraphrase experiments, stratify by problem difficulty, and extend cross-prover transfer. The mechanistic attribution framework demonstrated here---controlled ablation with falsification thresholds---provides a template for rigorous evaluation beyond theorem proving.

We began with a statistical gap and a hybrid puzzle. We end with a mechanistic answer: natural language is not incidental to LLM theorem proving---it's the dominant mechanism. This shifts the question from ''how do we make LLMs better provers?'' to ''how do we design systems that optimally exploit linguistic understanding while preserving formal reliability?'' The answer lies in hybrid architectures informed by mechanistic attribution.
